\centerline{\bf \Huge ДЗ. Олимпиадное движение} 


\problem{\textbf{1}} АА начал следить за механическими настенными часами в 8:30:52, а закончил в 13:40:57. Пусть X - это количество моментов, которое насчитал АА когда минутная стрелка и часовая лежали на одной прямой. Пусть Y - это количество моментов, которое насчитал АА когда минутная стрелка и секундная лежали на одной прямой. Чему равно значение X+Y?

\problem{\textbf{2}}  Кабинки горнолыжного подъёмника занумерованы подряд числами от 1 до 2026. Начальник сел в кабинку №67 подъёмника
у подножия горы и в какой-то момент заметил, что он поравнялся с движущейся вниз кабинкой №993, а через 52 секунды его кабинка поравнялась
с кабинкой №992. Через какое время Начальник прибудет на вершину горы?

